\begin{frame}{2층 신경망으로 XOR 풀기 + 역사}
  \begin{itemize}
    \item 은닉 각 유닛이 $z = \max(0, w_1x_1 + w_2x_2 + b)$ ---
      \textbf{한 개의 직선 경계}.
    \item 방향이 다른 두 유닛(``$x_1{+}x_2$ 큰 쪽''과
      ``$x_1{-}x_2$ 큰 쪽'')의 \textbf{두 직선}이 2차원을
      \textbf{4개 영역}으로 쪼갬.
    \item 출력층이 각 영역에 다른 가중치를 두어 \textbf{파란 두 점에만
      높은 값}. (실제 학습엔 \textbf{4개 이상}의 유닛 필요 ---
      2유닛은 잘 알려진 로컬 미니멈에 빠지기 쉬움.)
  \end{itemize}
  \vskip0.4em
  \begin{block}{역사적 맥락}
    XOR은 1969년 민스키$\cdot$파핏의 저서 \textit{Perceptrons}에서
    ``단층 퍼셉트론으로 풀 수 없는 가장 간단한 문제''로 지적되며
    \textbf{신경망 연구 10년 정지}의 결정적 원인.
    1986년 \textbf{역전파}(Rumelhart-Hinton-Williams)로 해소 ---
    2층 이상은 XOR을 풀고, 그 그래디언트를 계산하는 방법이 있음.
  \end{block}
\end{frame}
